\chapter{结论与展望}
\label{chap:concl}

\section{工作总结}

本文在 xv6-riscv 上设计并实现了并行于根文件系统的 \textbf{PMFS v3}：固定挂载点 \texttt{/pm} 上的受限文件语义、COW 元数据与不可变 KV run、双槽 checkpoint 提交、行级 flush 请求与可选 Zicbom，以及面向 QEMU 的实验性独立 NVDIMM 补丁路径。系统以进程重启切点 oracle、有界规格负控与轨迹采样验收提交行为，并以对齐工作量基准报告相对原生 \texttt{fs+log} 的软件路径延迟/吞吐量级差。

主要结论：

\begin{enumerate}[leftmargin=2em]
  \item \textbf{方法学：}不可替代点是可复现的验收资产与诚实假设清单，而非又一个双槽结构或完整生产 FS。
  \item \textbf{系统：}成功 \texttt{fsync} 后可读回完整版本；空白/损坏拒绝挂载；PM 路径不经块日志栈。
  \item \textbf{实验：}架构面展示量级优势但持久契约不等价；公平面展示软件请求效应，二者不可混读。
\end{enumerate}

对照原题四要件：NVDIMM 映射在\textbf{实验补丁路径}有条件达标；去块布局实质达标；刷写一致性在\textbf{声明假设}下达标而物理掉电未证；架构级性能叙事在软件路径上成立。

\section{不足}

真机持久域校准、上游通用设备发现、完整 POSIX、可扩展并发、近满可持续 GC、以及论文正式模板化定稿仍不足。旧 PDF/历史不等量基准不得作为新结论。

\section{展望}

\begin{itemize}[leftmargin=2em]
  \item 与导师书面确认设备模型与并行 \texttt{/pm} 产品形态；
  \item 将实际 store/flush 轨迹更紧地接入状态空间检查；
  \item 在保持尖刺的前提下扩展可测语义边界与论文图表重绘。
\end{itemize}

答辩一句话：\textbf{在教学 OS 上交付可经普通文件 API 访问的 PM 文件子系统，并以可脚本化故障验收与诚实假设证明提交协议；相对原生块日志展示架构级软件路径优势，但不声称无条件物理掉电或真机持久加速。}
